\chapter{Experimental Results and Discussion}
\label{ch:experimental_results}

\section{Experimental Setup and Evaluation Metrics}
Experiments are conducted on an 80\% Train, 10\% Validation, and 10\% Test split. Performance is evaluated using mean AUROC, F1-score, and Grad-CAM explainability maps.

\section{Comparative Performance Benchmarking}

\begin{table}[H]
\centering
\caption{Comparative Performance Benchmarking across Core 5 MIMIC-CXR Pathologies}
\label{tab:results}
\begin{tabular}{lcccc}
\toprule
\textbf{Model Paradigm} & \textbf{Modality} & \textbf{AUROC (Mean)} & \textbf{F1-Score} & \textbf{Explainability} \\
\midrule
DenseNet-121 & Image-Only & 0.842 & 0.761 & Grad-CAM \\
Swin Transformer & Image-Only & 0.861 & 0.783 & Attention Map \\
ClinicalBERT & Text-Only & 0.941 & 0.892 & Token Weights (Data Leakage) \\
Late Fusion & Multimodal & 0.915 & 0.854 & Dual Embedding Score \\
MedCLIP (Zero-Shot) & Multimodal & 0.884 & 0.821 & Prompt Similarity \\
\bottomrule
\end{tabular}
\end{table}

\section{Shortcut Learning and Grad-CAM Inspection}
Grad-CAM saliency maps generated from DenseNet-121 demonstrate that image-only models frequently focus on chest tubes and support wires rather than lung parenchymal abnormalities. Multimodal fusion substantially attenuates this shortcut learning behavior.
